\subsubsection{Congruencia ternaria de la traza y optimalidad del exponente}
\label{mg09:congruencia}

La identidad de nivel tres permite comparar todos los coeficientes
de la traza construida con los de su realización modular. El dominio
adecuado para esta consecuencia es el anillo de series formales: cada
coeficiente depende de un número finito de factores del producto,
mientras que el enunciado conserva todos los grados.

\begin{lemma}[Unidad integral del producto de orden tres]
La serie
\begin{equation}
 u(q)=q\,t_3(q)=\prod_{n\geq1}(1+q^n+q^{2n})^{-12}
 \label{mg09:unidad}
\end{equation}
pertenece a \(1+q\mathbb Z[[q]]\), y su inversa pertenece al
mismo conjunto.
\end{lemma}

\begin{proof}
Cada polinomio \(1+q^n+q^{2n}\) tiene término constante uno.
Para cualquier serie \(v(q)=1+\sum_{r\geq1}v_rq^r\) de
coeficientes enteros, la ecuación \(v(q)w(q)=1\) determina
recursivamente
\[
 w_0=1,\qquad w_r=-\sum_{i=1}^{r}v_iw_{r-i}\quad(r\geq1).
\]
Esta recurrencia prueba existencia, unicidad e integralidad de la
inversa. Aplicada a cada factor, también prueba la integralidad de su
potencia inversa doce. Además,
\((1+q^n+q^{2n})^{-12}\in1+q^n\mathbb Z[[q]]\).
Para calcular un coeficiente de grado a lo sumo \(r\), bastan los
factores con \(n\leq r\); los productos parciales estabilizan
coeficiente a coeficiente. Por tanto, \eqref{mg09:unidad} define
una serie integral con término constante uno. La misma recurrencia,
aplicada ahora a \(u\), da \(u^{-1}\in1+q\mathbb Z[[q]]\).
\end{proof}

\begin{proposition}[Divisibilidad uniforme exacta]
Para la traza \(t_3\) y la serie \(J=j-744\) relacionadas por
\eqref{km:j-nivel3}, se tiene
\begin{equation}
 J-(t_3+12)\in3^9q\mathbb Z[[q]].
 \label{mg09:congruencia-exacta}
\end{equation}
El exponente uniforme de divisibilidad por tres es exactamente nueve:
\begin{equation}
 [q]\bigl(J-t_3-12\bigr)=10\cdot3^9=196830.
 \label{mg09:optimalidad}
\end{equation}
\end{proposition}

\begin{proof}
Como \(t_3=q^{-1}u\), el lema da
\(t_3^{-r}=q^ru^{-r}\in q^r\mathbb Z[[q]]\) para
\(r=1,2,3\). La identidad \eqref{km:j-nivel3} se transforma en
\begin{equation}
 J-t_3-12
 =3^9q\left(10u^{-1}+4\cdot3^5q\,u^{-2}
                         +3^9q^2u^{-3}\right).
 \label{mg09:factorizacion}
\end{equation}
Todos los términos del paréntesis son series enteras; esto prueba
\eqref{mg09:congruencia-exacta}. Su término constante es \(10\),
pues \(u(0)=u^{-1}(0)=1\), y los dos últimos términos contienen
\(q\) y \(q^2\). Se obtiene \eqref{mg09:optimalidad}. Como
\(3\nmid10\), el coeficiente de \(q\) tiene valoración
\(3\)-ádica nueve, lo que fija la optimalidad del exponente uniforme.
\end{proof}

En particular, \([q^r]J\equiv[q^r]t_3\pmod{3^9}\) para todo
\(r\geq1\). El cuantificador procede de la identidad formal y de
la inversión integral demostrada en el lema. Una verificación hasta
cualquier grado finito es una especialización de esta relación.
La traza sigue procediendo del automorfismo de orden tres del retículo
construido desde la incidencia HMT; la uniformización actúa después
sobre esa traza. La congruencia conserva ese orden y determina una
relación exacta entre sus dos lecturas graduadas.
